\section{Introduction}
\label{sec:intro}

EEG foundation models tokenize a recording before they model it, and the
choice of what a token carries decides what the model can learn cheaply and
what it can learn only by accident. Across the current generation the
choice varies in how much of the signal survives tokenization --- magnitude
spectrograms \citep{yang2023biot,pradeepkumar2026tfm}, raw patches
\citep{wang2025cbramod,wang2024eegpt}, per-bin spectral amplitude and phase
\citep{jiang2024labram} --- but not in one respect: every token describes
one frequency band, or one patch of the broadband waveform, on its own. No
token type represents the \emph{relation} between two bands.

That relation is not an exotic quantity. Phase--amplitude coupling, the
locking of a fast band's envelope to the phase of a slow rhythm, has a
standard estimator \citep{tort2010measuring}, a physiology documented for
two decades \citep{canolty2006high,canolty2010functional}, and direct
clinical relevance to precisely the tasks foundation models are evaluated
on: it is stronger in the seizure onset zone than in unaffected tissue,
interictally as well as ictally \citep{amiri2016phase}, and periodic
epileptiform discharges are, by definition, fast activity riding a slow
rhythm. A representation from which such a quantity could in principle be
decoded is not yet one that encodes it. Magnitude tokenizers make it hard
to reach, because re-deriving phase from overlapping magnitude frames is a
phase-retrieval problem \citep{griffin1984signal} that nothing in the
pipeline performs; waveform tokenizers keep it reachable but give the
encoder no incentive to extract it, and probing studies find that the
resulting embeddings under-represent exactly the oscillatory structure
such relations live in \citep{kommineni2026aperiodic}.

This paper puts the relation into the token and asks a question that turns
out to have a sharp answer: \emph{when does a model need it there?} We
build a tokenizer that decomposes each electrode into learned frequency
bands, measures the coupling between them per patch, and hands it to the
encoder as a first-class token with three guarantees (phase-reference
invariance, magnitude preservation, exact recoverability of the uncoupled
model). We then hold the tokenizer fixed and vary the encoder. An encoder
that attends along a frequency axis can compute cross-band relations for
itself, and for it the coupling token is redundant on eleven of twelve
corpora. A published foundation-model encoder that keeps its own spectral
branch behaves the same way. An encoder with \emph{no} cross-frequency
pathway --- no frequency axis, no spectral branch, bands meeting only
inside the spatial attention --- makes the token load-bearing: removing it
costs between $0.03$ and $0.14$ on every one of the four corpora on which
we measure it. Coupling is therefore worth carrying as \emph{content}
exactly when the encoder does not supply it as \emph{structure}, and the
model built on that principle is the smallest and, on seizure detection,
the strongest in our comparison.

\paragraph{Contributions.}
\begin{itemize}[leftmargin=*,itemsep=2pt,topsep=2pt]
\item \textbf{Cross-frequency coupling as token content.} Each band
  contributes a waveform token and an interaction token built from the
  band's measured coupling to slower bands; the interaction token is
  invariant to the phase reference of the analysis, preserves band energy
  exactly, and is gated so that the uncoupled model is an exactly
  reachable configuration (Section~\ref{sec:tokenizer}).
\item \textbf{An axis-free encoder designed for it, and the attribution
  that goes with it.} CroFreMo folds the band rows into spatial attention
  and has no frequency sub-layer; at 2.7M parameters it is competitive
  with 25--69M foundation models on twelve clinical corpora and ahead of
  the strongest reproduced baseline by $0.15$ and $0.09$ AUC-PR on the two
  seizure-detection corpora with adequate positives. With the encoder
  fixed, switching the coupling token off costs $0.136$ on CHB-MIT,
  $0.107$ on TUEV, $0.063$ on TUSZ and $0.027$ on IIIC
  (Section~\ref{sec:attribution}): the seizure margin is made of coupling
  content.
\item \textbf{Content versus structure, measured across three encoders.}
  The same token inside a tri-axial encoder with a frequency axis is
  decisive on one corpus and neutral on the rest; added to CBraMod's
  encoder next to its own spectral branch, it is attributable on
  event-classification corpora and harmful on seizure detection; inside
  the axis-free encoder it is necessary everywhere we look
  (Section~\ref{sec:content-structure}). The scope of a tokenizer claim is
  set by the encoder it is paired with, and we report it that way.
\item \textbf{A pretraining study reported to convergence rather than
  assumed.} We construct the masked-reconstruction objective that
  supervises coupling directly, evaluate it with lr-matched frozen probes,
  label-fraction curves and clean-transfer controls, and report the
  pretrained tri-axial model under one fixed finetuning recipe: it helps
  where labels are scarce and hurts where they are plentiful
  (Section~\ref{sec:pretrain-results}).
\end{itemize}

Two boundaries are stated up front. The axis-free model wins or ties on
the six corpora whose label is a paroxysmal event, a pattern, or a
recording-level abnormality, and loses to a pretrained or a larger
from-scratch model on sleep staging, cohort-level diagnosis, artifact
labelling and the smallest seizure corpus; those six cells are reported,
not averaged away. And the cells of the axis-free model are single-seed
at the time of writing, against three-seed baselines; the attribution
effects are several times the seed variance we measure on the same
corpora, the main-table margins on the small corpora are not, and the
text says which is which.
